\chapter{Formalization-Oriented Theorem Decomposition}
\label{app:perverse-formalization}

This appendix is not Lean code.  It is a theorem-boundary specification
designed so that later Lean files can mirror the mathematical dependency
graph without introducing a single project-specific axiom.

\section{Generic categorical layer}

The first reusable layer should not mention finite fields:

\begin{enumerate}
\item a bounded \(t\)-structure on a triangulated or stable category;
\item its heart, truncation triangles, and cohomology functors;
\item gluing of \(t\)-structures along recollement data;
\item exact and one-sided \(t\)-exact functors;
\item finite filtrations supplied by a truncation tower;
\item a duality equivalence exchanging the two halves of a self-dual
      \(t\)-structure;
\item abelian images and the abstract construction of an intermediate object
      as \(\operatorname{im}(j_!\to j_*)\).
\end{enumerate}

Theorems in this layer apply to coherent, constructible, and motivic
\(t\)-structures as well as to perverse sheaves.

\section{Dimension and support layer}

For a noetherian topological space equipped with a dimension function
\(\delta\), define predicates
\[
\begin{aligned}
\operatorname{PervLE}(K)&:\Longleftrightarrow
 H^a(i_x^*K)=0\ \text{for }a>-\delta(x),\\
\operatorname{PervGE}(K)&:\Longleftrightarrow
 H^a(i_x^!K)=0\ \text{for }a<-\delta(x).
\end{aligned}
\]
The terminal construction theorem should return a bounded \(t\)-structure
whose two classes are propositionally equal to these predicates.  Important
test cases are:
\[
 \operatorname{Perv}(\Spec k)=\operatorname{Shv}_c(\Spec k),\qquad
 \operatorname{IsLisse}(L)\Rightarrow L[d]\in\operatorname{Perv}(X)
\]
for \(X\) smooth and pure \(d\)-dimensional.

\section{Intermediate extension layer}

Given an open--closed recollement, expose:
\[
 {}^pj_!P={}^pH^0(j_!P),\qquad
 {}^pj_*P={}^pH^0(Rj_*P),
\]
and a natural transformation between them.  Define
\[
 \operatorname{middleExtension}(P)
 =
 \operatorname{image}({}^pj_!P\to{}^pj_*P).
\]
The principal theorem is an equivalence among:

\begin{enumerate}
\item being isomorphic to this image;
\item restricting to \(P\) and having no boundary quotient;
\item restricting to \(P\) and having no boundary subobject;
\item satisfying both boundary conditions.
\end{enumerate}

Transitivity and choice independence should be proved by this universal
property.  The intersection complex is then a structure consisting of a
smooth dense stratum, a lisse coefficient object, and the canonical
middle-extension isomorphisms under shrinking.  It must not store a
noncanonical chosen stratification in the public theorem statement.

\section{Weight layer}

Keep the following types distinct:

\begin{enumerate}
\item \texttt{PunctuallyPure iota w F} for a Weil sheaf;
\item \texttt{WeightsLE iota w K} for a complex;
\item \texttt{WeightsGE iota w K}, defined through duality;
\item \texttt{Pure iota w K} as their conjunction;
\item \texttt{UniversallyPure w K}, quantifying over all complex
      realizations.
\end{enumerate}

Prove the normalization lemmas
\[
\begin{aligned}
\operatorname{Pure}(w,K[n])&\leftrightarrow
 \operatorname{Pure}(w-n,K),\\
\operatorname{Pure}(w,K(m))&\leftrightarrow
 \operatorname{Pure}(w+2m,K),\\
\operatorname{Pure}(w,\D K)&\leftrightarrow
 \operatorname{Pure}(-w,K),
\end{aligned}
\]
and the perverse-cohomology bridge
\[
\operatorname{WeightsLE}(w,K)
\Longleftrightarrow
\forall a,\ \operatorname{WeightsLE}(w+a,\pH^aK).
\]
The smooth normalization theorem should literally implement
\[
\begin{gathered}
\operatorname{PunctuallyPure}(0,\Qbarl)\\
\Longrightarrow \operatorname{Pure}(d,\Qbarl[d]).
\end{gathered}
\]

\section{Fourier layer}

For a vector bundle \(E\to S\) of rank \(r\), package the pairing
\(E\times_SE^\vee\to\A^1\) and the Artin--Schreier kernel.  Define
\[
 \FT_\psi(K)=
 Rp_{2!}\bigl(p_1^*K\otimes\cL_\psi(\langle x,\xi\rangle)\bigr)[r].
\]
The reusable theorem suite is:

\begin{enumerate}
\item base change and the stalk formula;
\item Fourier inversion, with the sign map and Tate twist explicit;
\item compatibility with duality;
\item trace-function finite Fourier transform;
\item Plancherel;
\item perverse \(t\)-exactness.
\end{enumerate}

Weight preservation must not be part of this upstream suite.  It is a later
corollary of the direct-image theorem.

\section{Direct-image theorem layer}

The main theorem should first be stated in derived form:
\[
 \operatorname{WeightsLE}(w,K)
 \Longrightarrow
 \operatorname{WeightsLE}(w,Rf_!K).
\]
Its implementation dependencies should be separate lemmas for:

\begin{enumerate}
\item quasi-finite maps;
\item composition via a finite filtration;
\item open--closed devissage;
\item relative Noether normalization;
\item affine-line projection using the curve estimate;
\item proper-support base change at a geometric point;
\item production of mixedness via real envelopes.
\end{enumerate}

The perverse statement is a corollary, not a second proof:
\[
 \operatorname{WeightsLE}(w,P)
 \Rightarrow
 \operatorname{WeightsLE}(w+a,\pH^aRf_!P).
\]

\section{Proper purity and decomposition layer}

Implement the chain as separate declarations:
\[
\begin{aligned}
\texttt{rfStarWeightsGE}&:
 \operatorname{WeightsGE}(w,K)\to
 \operatorname{WeightsGE}(w,Rf_*K),\\
\texttt{properPushforwardPure}&:
 \operatorname{Proper}(f)\to\operatorname{Pure}(w,K)\to
 \operatorname{Pure}(w,Rf_*K),\\
\texttt{middleExtensionPure}&:
 \operatorname{Pure}(w,P)\to\operatorname{Pure}(w,j_{!*}P),\\
\texttt{intersectionCohomologyPure}&:
 \operatorname{Proper}(X)\to
 \operatorname{Pure}(w,\IC_X(L))\to
 \operatorname{Pure}(w+r,IH^r).
\end{aligned}
\]

Only after these should one add:

\begin{itemize}
\item geometric semisimplicity of pure perverse objects;
\item existence of a splitting of a pure complex;
\item the decomposition theorem for proper pushforward.
\end{itemize}

The splitting theorem should return a merely existing isomorphism after base
change.  It should not manufacture a canonical splitting or a theorem that
Frobenius is semisimple.

\section{Final assembly}

For a smooth proper \(d\)-fold, the key Lean term should visibly compose
\[
\begin{tikzcd}
\Qbarl[d]\arrow[r,"\text{proper purity}"]&
Ra_*\Qbarl[d]\arrow[r,"H^{i-d}"]&
H^i(X,\Qbarl).
\end{tikzcd}
\]
The final theorem then combines this result with three geometry-free modules:

\begin{enumerate}
\item alternating determinant--trace and the Euler product;
\item similitude pairings and the functional equation;
\item rational-series integrality, universal-weight separation, and
      \(\ell\)-independence.
\end{enumerate}

Every release theorem must pass an axiom audit.  Temporary capability records
may support parallel development, but no theorem in the release import cone
may depend on a project-specific axiom, a placeholder interface, or a
noncomputable choice whose independence has not been proved.
